% TOP_PROBLEM: {"id": 30006934, "title": "Piltz Divisor Problem for k-Fold Divisor Error Terms", "category_id": 1, "category": {"id": 1, "name": "number_theory", "display_name": "Number Theory", "description": "Properties of integers, prime numbers, Diophantine equations.", "slug": "number-theory", "order_index": 1, "created_at": "2026-07-31T15:26:25.670Z"}, "status": "open"}
% FIRST_POSTED: 2026-09-25
% ENTRY_KIND: statement_only
% !TeX program = lualatex
% Definition and sources only; this file is not an LLM solution attempt.
\documentclass[11pt]{article}
\usepackage[margin=25mm]{geometry}
\usepackage{fontspec}
\setmainfont{Latin Modern Roman}
\usepackage{amsmath,amssymb,mathtools}
\usepackage{unicode-math}
\setmathfont{Latin Modern Math}
\usepackage{xurl,hyperref}
\hypersetup{colorlinks=true,urlcolor=blue}
\setcounter{secnumdepth}{0}
\setlength{\parindent}{0pt}
\setlength{\parskip}{5pt}
\setlength{\emergencystretch}{3em}
\begin{document}
\section*{\texorpdfstring{Piltz Divisor Problem for k-Fold Divisor Error Terms}{Piltz Divisor Problem for k-Fold Divisor Error Terms}}\label{30006934}
\subsection{Definitions and mathematical statement}
For $k\ge3$, let $d_k(n)$ count ordered $k$-tuples of positive integers with product $n$. Define the error in the summatory function by
\[
 \Delta_k(x)=\sum_{n\le x}d_k(n)
 -\operatorname*{Res}_{s=1}\frac{\zeta(s)^k x^s}{s},\qquad x\ge1.
\]
The residue is the complete polynomial-in-$\log x$ main term multiplied by $x$, not just its highest-degree term. Define
\[
 \alpha_k=\inf\{a:\ \text{for every }{\varepsilon}>0,
 \ \Delta_k(x)=O_{k,{\varepsilon}}(x^{a+{\varepsilon}})\}.
\]
Determine $\alpha_k$ for every $k\ge3$. Constants may depend on $k$ and ${\varepsilon}$, but not on $x$. This is a pointwise error-exponent problem rather than a mean-square or short-interval average problem.
\par\smallskip\textit{Formulation and concept references:} \href{https://arxiv.org/html/2303.05028}{[S1]}.

\subsection{Short English statement}
What is the best possible pointwise error exponent when counting ordered factorizations into a fixed number of positive factors?
\subsection{Sources}
\begin{itemize}
\item[S1]Bellotti and Yang, On the generalised Dirichlet divisor problem (2023 preprint; BLMS 2024).
\url{https://arxiv.org/html/2303.05028}.
\textit{Formulation source.}
\item[S2]Status source for Piltz Divisor Problem for k-Fold Divisor Error Terms.
\url{https://academic.oup.com/imrn/article/2024/20/13191/7755133}.
\textit{Further reference; not a proof endorsement.}
\item[S3]Status source for Piltz Divisor Problem for k-Fold Divisor Error Terms.
\url{https://arxiv.org/html/2506.22587v1}.
\textit{Further reference; not a proof endorsement.}
\end{itemize}
\end{document}
